% Standalone note. From the repository root, compile with:
% latexmk -pdf -outdir=.vscode/latex-build/notes notes/corruption_kd_lower_bound.tex
% Copy the resulting PDF to notes/corruption_kd_lower_bound.pdf if desired.
\documentclass[11pt]{article}

\usepackage[T1]{fontenc}
\usepackage{lmodern}
\usepackage[margin=1in]{geometry}
\usepackage{amsmath,amssymb,amsthm}
\usepackage{hyperref}
\hypersetup{pdftitle={A matching lower bound for corrupted verifier feedback}}

\newcommand{\X}{\mathcal X}
\newcommand{\C}{\mathcal C}
\newcommand{\V}{\mathcal V}
\newcommand{\E}{\mathbb E}
\newcommand{\ind}{\mathbf 1}
\newtheorem{theorem}{Theorem}
\newtheorem{lemma}[theorem]{Lemma}

\title{A matching lower bound for corrupted verifier feedback}
\author{}
\date{}

\begin{document}
\maketitle

We construct a verifier class whose optimal deterministic realizable mistake
bound is exactly $d$, but for which at most $k$ corrupted responses can force
every randomized cooperative learner to make at least $d(k+1)/4$ expected
mistakes. The same class works for every $k$. The verifier is fixed before
the interaction, while the corruption decisions may depend on the queried
certificates.

The idea is to hide one certificate value among $4d$ possibilities. Each
query tests at most two of them. Without corruption, testing every value
once determines the true label. With corruption, the first $k$ positive
responses can be hidden. We show that many possible hidden values remain
indistinguishable from a verifier whose output is always zero.

\section{Protocol and verifier class}\label{sec:class}

Fix an integer $d\ge1$. Let $\X=\{x_0\}$ and $\C=[4d]=\{1,\ldots,4d\}$.
The input is $x_0$ on every round. We first consider debates of length one.
A cooperative learner predicts a label $\hat y_t\in\{0,1\}$ and chooses
certificates $p_t,q_t\in\C$. It then observes one verifier response
$\tilde y_t$. The true label is not revealed separately, and no additional
verifier queries are allowed. Both the prediction and the certificates may
depend on the earlier interaction and the learner's private randomness.

For a fixed verifier $v^*$, let $h_{v^*}$ denote its validated hypothesis.
We count prediction mistakes by
\[
M(T)=\sum_{t=1}^T\ind\{\hat y_t\ne h_{v^*}(x_0)\}.
\]
A round is corrupted when $\tilde y_t\ne v^*(x_0,p_t,q_t)$. The number of
corrupted rounds is at most $k$ on every run.

For $b\in\{0,1\}$, let $h_b(x_0)=b$. Define the verifiers
\[
v_0(x_0,p,q)=0,
\qquad
v_j(x_0,p,q)=\ind\{p=j\text{ or }q=j\},\quad j\in[4d],
\]
and let $\V_d=\{v_0\}\cup\{v_j:j\in[4d]\}$.
Think of each verifier as a $4d\times4d$ table, with the first certificate
indexing its row and the second indexing its column. The table of $v_0$ is
all zero. The table of $v_j$ is zero except that row $j$ and column $j$ are
all one.

These are valid verifiers. The verifier $v_0$ validates $h_0$. For $v_j$,
the proposer can force output $1$ by choosing $p=j$. No choice of $p$ can
force output $0$, because the responder can always choose $q=j$ and obtain
output $1$. Thus every $v_j$ validates $h_1$, and the supported hypothesis
class is precisely $\{h_0,h_1\}$.

\section{The deterministic realizable mistake bound}\label{sec:clean}

We can determine the true label by querying disjoint pairs until every
certificate value has been tested. Alternating the predictions during this
search limits the mistakes for either possible true label.

\begin{lemma}\label{lem:clean}
The optimal deterministic realizable mistake bound for $\V_d$ is exactly $d$.
\end{lemma}
\begin{proof}
For the upper bound, query the pairs
\[
(1,2),(3,4),\ldots,(4d-1,4d),
\]
in this order, while alternating predictions $0,1,0,1,\ldots$. If a response
is $1$, the verifier cannot be $v_0$, so predict $1$ on every later round.
If all $2d$ responses are zero, the verifier must be $v_0$, so predict $0$
on every later round.

When the true label is $0$, all responses are zero and exactly $d$ of the
first $2d$ predictions are wrong. When the true label is $1$, a positive
response occurs within those $2d$ rounds. Only predictions of $0$ before
or on that round can be wrong, and there are at most $d$ of them. Thus the
learner makes at most $d$ mistakes for every horizon.

For the lower bound, fix any deterministic learner and follow its first
$2d-1$ rounds when every response is zero. Its queries use at most $4d-2$
distinct certificate values, so there is an index $j\in[4d]$ that never
appears. Both $v_0$ and $v_j$ give exactly this response sequence. Among the
$2d-1$ predictions, one label appears at least $d$ times. If that label is
$1$, choose $v^*=v_0$; if it is $0$, choose $v^*=v_j$. In either case, the
learner makes at least $d$ mistakes. Since the learner is deterministic,
this entire response sequence and the chosen verifier can be determined
before the actual interaction.
\end{proof}

By the cooperative characterization in the manuscript, the debate
Littlestone dimension of this class is therefore $d$.

\section{The lower bound with corrupted feedback}\label{sec:corrupt}

We now hide the first $k$ positive responses of a verifier $v_j$. As long as
the learner has queried the hidden value $j$ at most $k$ times, all observed
responses are zero. This is exactly what the learner would observe from
$v_0$, whose true label is opposite. The following argument counts how many
hidden values can still be concealed after a given number of queries.

\begin{theorem}\label{thm:corrupt}
For every integer $k\ge0$, every $T\ge d(k+1)$, and every randomized
cooperative learner for $\V_d$, there is a fixed verifier $v^*\in\V_d$ and
an adaptive corruption strategy using at most $k$ corruptions on every run
such that $\E[M(T)]\ge d(k+1)/4$. The expectation is over the learner's
randomness, and the verifier is chosen before that randomness is observed.
\end{theorem}
\begin{proof}
Let $m=d(k+1)$. We first give a distribution over verifiers, chosen
independently of the learner's random bits before interaction. Choose $v_0$
with probability $1/2$. With the remaining probability $1/2$, choose an
index $J$ uniformly from $[4d]$ and use $v_J$.

Under $v_0$, leave every response unchanged. Under $v_j$ for $j\in[4d]$,
replace its first $k$ responses equal to $1$ by $0$, and leave all other
responses unchanged. Apply this rule only during the first $m$ rounds;
afterward, leave every response unchanged. This strategy uses at most $k$
corruptions on every run. It changes only the observed responses, not the
fixed verifier or its validated hypothesis.

Fix a realization $r$ of all the learner's random bits. Run the resulting
deterministic policy hypothetically for $m$ rounds with all feedback equal
to zero. Write $(p_t^{(r)},q_t^{(r)})$ for its certificate pair and
$\hat y_t^{(r)}$ for its prediction on round $t$ of this hypothetical run.
For each $j\in[4d]$, let
\[
N_j(r)=\sum_{t=1}^m
\ind\{p_t^{(r)}=j\text{ or }q_t^{(r)}=j\}.
\]
This counts the rounds in which the query tests value $j$. Each round tests
at most two distinct values, so $\sum_{j=1}^{4d}N_j(r)\le2m$.
Consequently,
\[
\bigl|\{j:N_j(r)\ge k+1\}\bigr|
\le\frac{2m}{k+1}=2d.
\]
Define $G(r)=\{j\in[4d]:N_j(r)\le k\}$. We have $|G(r)|\ge2d$.

For every $j\in G(r)$, the actual corrupted interaction with $v_j$ follows
the hypothetical all-zero run for all $m$ rounds. To see this, proceed
inductively. As long as earlier feedback agrees, the learner chooses the
same certificates as in the hypothetical run. A query not containing $j$
has raw response zero. A query containing $j$ has raw response one, but it
is among at most $k$ such queries in this run, so the corruption strategy
changes it to zero. The transcripts therefore continue to agree.

Let $a(r)=\sum_{t=1}^m\ind\{\hat y_t^{(r)}=1\}$, the number of predictions
of $1$ on the hypothetical run. Also let $M_r(v;m)$ denote the actual
mistake count in the first $m$ rounds with verifier $v$, random bits $r$,
and the corruption strategy just defined. Under $v_0$, the true label is
$0$ and $M_r(v_0;m)=a(r)$. For each $j\in G(r)$, the predictions are the
same but the true label is $1$, so $M_r(v_j;m)=m-a(r)$.

Averaging over the verifier distribution while keeping $r$ fixed gives
\begin{align*}
\frac12 M_r(v_0;m)+\frac{1}{8d}\sum_{j=1}^{4d}M_r(v_j;m)
&\ge\frac12 a(r)+\frac{|G(r)|}{8d}\bigl(m-a(r)\bigr)\\
&\ge\frac12 a(r)+\frac14\bigl(m-a(r)\bigr)\\
&\ge\frac m4.
\end{align*}
This holds for every realization $r$. Averaging over the learner's random
bits shows that the expected mistake count under our distribution of fixed
verifiers is at least $m/4$. Hence at least one fixed verifier in its
support has expected mistake count at least $m/4$, where this last
expectation is over the learner's randomness alone. The corruption rule
for that verifier is already fixed. Finally, $M(T)\ge M(m)$ because
$T\ge m$, which proves the claim.
\end{proof}

The corruption budget in Theorem~\ref{thm:corrupt} holds on every run; it
is not merely a bound on the expected number of corruptions. The verifier
is chosen in advance, but the corrupted rounds depend on the learner's
queries. The proof does not require a schedule of corrupted rounds chosen
in advance.

\section{Longer debates}\label{sec:longer}

The same construction works for every debate length $n\ge1$. Given the
verifiers in Section~\ref{sec:class}, define
\[
v_j^{(n)}(x_0,c_1,\ldots,c_{2n})=v_j(x_0,c_1,c_2),
\qquad j\in\{0,1,\ldots,4d\},
\]
and let $\V_{d,n}=\{v_j^{(n)}:j\in\{0,1,\ldots,4d\}\}$. The later
certificates do not affect the response. For $j\in[4d]$, choosing $c_1=j$ still forces
output $1$ under $v_j^{(n)}$, and choosing $c_2=j$ still prevents the
proposer from forcing output $0$. Thus the validated hypotheses are
unchanged. Each round still tests at most two values and reveals just one
bit, so Lemma~\ref{lem:clean} and Theorem~\ref{thm:corrupt} apply without
changing either bound.

\end{document}
